\section{Introduction}

A recent multilingual evaluation reported an awkward pattern: every tested model was above chance, yet none beat the benchmark's majority-class baseline \citep{arcon2026metalinguistic}. The authors found the flip because they happened to print both reference lines. This is no curiosity of one newly constructed benchmark: a systematic construct-validity review of 445 language-model benchmarks offers 27 actionable checklist items, none of which asks authors to report a null, a chance level, or a constant predictor \citep{bean2025measuring}. The missing item is operational: \emph{before comparing model arms, test whether the reported construct clears its own input-blind floor}.

For a letter-valued multiple-choice construct, the floor is the most accurate constant letter under the benchmark's empirical gold marginal; for content likelihood, an analogous input-blind heuristic can exploit option length. These references need not equal nominal chance, and more importantly an arm can look ``above chance'' while carrying no more usable signal through the measured interface than a constant predictor --- especially under structural damage, where option-token selection bias \citep{zheng2024selectors} can dominate the argmax while content knowledge remains partly available.

We study null calibration as a \emph{measurement protocol}, not a new multiple-choice interface or debiasing algorithm. Majority-class baselines for MCQA already exist \citep{balepur2024artifacts}; letter-versus-content formulations are established \citep{gu2025olmes,cho2026choices}; constant outputs are known to game automatic judges \citep{zheng2025cheating}. Our contribution is to make the construct-appropriate input-blind null a pre-comparison gate, quantify the degrees of freedom hidden inside it, and separate two questions needing different references:

\begin{enumerate}[leftmargin=*,nosep]
    \item \textbf{Is the interface valid for arm comparison?} Use an arm-independent best-constant floor.
    \item \textbf{Does this particular arm carry item-level information?} Use an arm-conditional permutation null that preserves the arm's prediction marginal.
\end{enumerate}

The distinction prevents two opposite mistakes: chance can credit a literal constant emitter as competent, while the best-constant floor can label two equally empty emitters differently solely for collapsing onto different letters. The first motivates our headline floor; the second motivates read-out v2, for which we prove constant-collapse invariance algebraically.

Our evidence uses ten target constructs plus Winogrande as a negative control, four base-model families, letter and content likelihood interfaces, and structurally damaged arms. Our findings:

\begin{itemize}[leftmargin=*,nosep]
    \item The null is not ``chance,'' and the floor is itself an estimate needing its own calibration against a null the construct can \emph{realise} --- a test we initially failed on our largest benchmark, where uniform-over-ten-letters assigns gold letters 17\% of MMLU-Pro items cannot have. BoolQ's best constant is 0.6217, not 0.50; MMLU-Pro's floor is 0.116606, or 1.1661$\times$ naive ten-way chance. We report both gaps and ratios because either alone distorts comparisons across option counts. Against a null restricted to each item's own legal letters (Table~\ref{tab:nulls}), only MMLU and BoolQ have floors it could not plausibly produce ($p<10^{-5}$), and both have a genuinely fixed option count; the six others, MMLU-Pro included, lie inside the estimator's sampling noise ($p=0.083$--$0.853$). We therefore rest the quantitative claim on MMLU and BoolQ, reading the rest as showing only that chance is the wrong reference to \emph{report}, not that the floor is far from it. MMLU-Pro at $p=0.083$ is unresolved, not settled; Section~\ref{sec:legality-aware-null} says which rows moved and how.
    \item Content floors are not dataset constants unless their tie convention, length unit, and tokenizer are fixed: on MMLU-Pro, one tokenizer and one dataset yield floors from 0.125914 to 0.532164 across tie conventions, the latter exceeding the intact model's content score by 32.5 points.
    \item At MMLU-scale power, 17 designated damaged cells across four families exhibit wrong-null readings: 10/12 non-OLMo cells exceed item-averaged chance (12/12 exceed naive 0.10), yet only 1/12 clears the floor. The honest statement is 15/17 at or below the arm-independent floor, the two exceptions differing in kind --- \texttt{qwen3\_8b\_base/k14} is statistically real but materially negligible, while OLMo-2 \texttt{shortgpt16} retains 16 of 32 layers and clears its floor by $5.2\times$ its own half-width. Off MMLU the same set gives 9/85 above floor against 45/85 above chance, all nine healed arms retaining at least 14 layers; the accompanying power table is essential.
    \item The permutation null makes every pure constant emitter score exactly zero, re-sorts rather than shrinks the 27-cell read-out, and reveals intact Llama-2-7B as itself a weak MMLU-Pro anchor.
    \item Full-fp32 evaluation removes 100\% of bf16 exact ties and changes 2,532 letter decisions without recovering accuracy, so the failure is not repaired by numerical precision.
\end{itemize}

The floor test is deliberately one-directional: a failure disqualifies the score as evidence for an arm comparison, but a pass does not certify that the benchmark measures the intended capability. \citet{feng2019misleading} construct a dataset where a partial-input baseline is at chance while artifacts remain exploitable; clearing a null is thus \emph{necessary, not sufficient}.
